% BHU PYQ -- Manually transcribed & error-corrected
\documentclass[11pt,a4paper]{article}

\usepackage[a4paper,top=2.5cm,bottom=2.5cm,left=2.8cm,right=2.8cm,headheight=14pt]{geometry}
\usepackage{amsmath,amssymb}
\usepackage[T1]{fontenc}
\usepackage[utf8]{inputenc}
\usepackage{lmodern}
\usepackage{microtype}
\usepackage{fancyhdr}
\usepackage{enumitem}
\usepackage{hyperref}
\usepackage{lastpage}
\usepackage{parskip}

\hypersetup{colorlinks=true,urlcolor=black,linkcolor=black,
  pdftitle={BPE-602: Elementary Nanoscience and Exotic Materials -- BHU PYQ},pdfauthor={Banaras Hindu University}}

\pagestyle{fancy}\fancyhf{}
\renewcommand{\headrulewidth}{0.4pt}\renewcommand{\footrulewidth}{0.4pt}
\fancyhead[L]{\small\textit{Banaras Hindu University}}
\fancyhead[R]{\small\textit{BPE-602: Elementary Nanoscience and Exotic Materials}}
\fancyfoot[C]{\small Page \thepage\ of \pageref{LastPage}}

\newlist{parts}{enumerate}{2}
\setlist[parts,1]{label=(\alph*),leftmargin=*,itemsep=5pt,topsep=3pt}
\setlist[parts,2]{label=(\roman*),leftmargin=*,itemsep=3pt,topsep=2pt}
\newcommand{\pts}[1]{\hfill\textit{[#1]}}

\begin{document}

\begin{center}
  {\large\bfseries BANARAS HINDU UNIVERSITY}\\[2pt]
  {\normalsize B.Sc. (Hons.) Semester VI Examination, 2013-14}\\[6pt]
  \rule{0.8\linewidth}{0.5pt}\\[6pt]
  {\large\bfseries Paper No. BPE-602: Elementary Nanoscience and Exotic Materials}\\[4pt]
  {\normalsize Time: 3 Hours\hfill Maximum Marks: 70}
\end{center}
\rule{\linewidth}{0.4pt}
\smallskip
\noindent\textit{Instructions: Answer all questions. Symbols have their usual meanings.}
\bigskip

\noindent\textbf{Question 1.} Answer any \textbf{four} parts out of the following: \pts{4 $\times$ 3.5 = 14}
\begin{parts}
  \item Discuss the Fibonacci sequences and their relevance to quasi-crystals.
  \item Explain briefly the main features of Colossal Magnetoresistance (CMR).
  \item How is focusing of electrons performed in a transmission electron microscope (TEM)? Name the two lenses used in a TEM.
  \item In which characterization technique are secondary electrons exploited? Discuss the interaction of electrons with thin samples.
  \item Define Quantum dot, quantum wire, and quantum well with suitable examples.
\end{parts}

\medskip\rule{\linewidth}{0.2pt}

\noindent\textbf{Question 2.}
\begin{parts}
  \item What do you understand by tunneling? Discuss the salient features of scanning probe microscope (SPM) and write the construction and working principle of scanning tunneling microscope (STM). \pts{4 + 10 = 14}
\end{parts}

\medskip\rule{\linewidth}{0.2pt}

\noindent\textbf{Question 3.}
\begin{parts}
  \item What are the essential differences between crystals and quasi-crystals? Compare them in a tabulated form. Why does one need a six-index system in order to index quasi-crystalline diffraction patterns? \pts{14}
\end{parts}

\medskip\rule{\linewidth}{0.2pt}

\noindent\textbf{Question 4.}
\begin{parts}
  \item Discuss the structure, properties, and applications of carbon nanotubes (CNTs). Discuss in detail the synthesis of carbon-based nanomaterials using DC arc discharge techniques. \pts{9}
  \item Write a short note on the arc-discharge method for synthesizing single-walled carbon nanotubes (SWNTs). \pts{5}
\end{parts}
\hfill \textbf{OR}
\begin{parts}
  \item Discuss with an appropriate diagram the mechanism of Giant Magnetoresistance (GMR). \pts{10}
  \item Explain the spin valve effect. \pts{4}
\end{parts}

\medskip\rule{\linewidth}{0.2pt}

\noindent\textbf{Question 5.}
\begin{parts}
  \item What are superconductors? How are they different from perfect conductors? Explain the Meissner effect and supercurrents in superconductors. \pts{14}
\end{parts}
\hfill \textbf{OR}
\begin{parts}
  \item Discuss the salient features of High-Temperature Superconductors (HTS). Draw a neat diagram of the structure of Y-123 ($\text{YBa}_2\text{Cu}_3\text{O}_{7-x}$) and point out the $\text{Cu--O}$ plane in which charge carriers exist. \pts{14}
\end{parts}

\vfill
\rule{\linewidth}{0.4pt}
\begin{center}\small
  \href{https://bhu-exam.vercel.app/}{Click here to visit our website}\\[4pt]
  \href{https://me-aryan.vercel.app/}{Click here to visit our contributor}\\[4pt]
  \href{https://quashan-me.vercel.app/}{Click here to visit our contributor}
\end{center}

\end{document}
